\section{Empirical Methodology}\label{sec:method}

\subsection{Two Structurally Distinct Estimators of the Aggregate Level}\label{sec:method-estimators}

The first estimator is a 10,000-household Agent-Based Model (ABM), specified in Appendix~\ref{sec:method-abm}. It draws a synthetic population calibrated to aggregate FRED medians. It then tests whether a utility-maximizing household decision rule, rational or behaviorally extended, can account for the shortfall; Section~\ref{sec:abm} reports its results. The second is a hazard framework, specified in Section~\ref{sec:hazard}. It abandons the household utility function entirely and models loan-level survival directly on real Freddie Mac loan-level data. It runs along two paths. Path A is a stratum-month discrete-time hazard model, estimated by Poisson pseudo-maximum likelihood on the full 2017--2021 origination universe aggregated to stratum-months. Path B is a loan-level microsimulation with competing prepayment and default risks, run on a stratified 75,000-loan sample drawn from that universe.

I evaluate both estimators against the same \$764.7 billion empirical benchmark (Section~\ref{sec:method-benchmark}) and the same June 2022--November 2025 QT window. Their results are therefore directly comparable, even though they draw on different data and different behavioral assumptions. They differ in modeling paradigm --- household choice versus loan-level survival --- and, on the hazard side, in the use of real loan-level data. Population size differs too: 10,000 households against Path B's 75,000 simulated loans, with Path A estimated on the full origination universe. Section~\ref{sec:robustness} finds that this size difference does not drive the gap in their explanatory power. The data-source difference, though, remains an unresolved confound that qualifies the paradigm interpretation throughout. Figure~\ref{fig:architecture} maps the pipeline.

\begin{figure}[H]
\centering
\includegraphics[width=\textwidth]{fig7_architecture}
\caption{Pipeline architecture. Data sources (FRED macro series, NY Fed SOMA holdings, Freddie Mac loan-level performance) feed three estimators (the synthetic-population ABM and hazard Paths A and B), whose simulated roll-off paths are scored through one shared macro-accounting layer against the \$764.7 billion empirical benchmark of Section~\ref{sec:method-benchmark}. The two robustness cross-links are the cross-design test (real Freddie covariates into the ABM, Section~\ref{sec:robustness-crossdesign}) and the symmetric companion test (synthetic population into the hazard framework, Section~\ref{sec:robustness-synthesis}).}
\label{fig:architecture}
\end{figure}

The agent-based model's specification, and the structural extensions whose grounding Section~\ref{sec:abm}'s results reference, are given in full in Appendix~\ref{sec:method-abm} and Appendix~\ref{sec:method-frictions}.

\subsection{Empirical Validation Benchmark}\label{sec:method-benchmark}

I build the empirical shortfall benchmark from the New York Fed's SOMA current face value series. It nets realized roll-off against the actual phased redemption caps --- \$17.5 billion per month June--August 2022, \$35 billion per month from September 2022 --- strictly within the active QT window, June 2022--November 2025. Accumulation is net rather than shortfall-months-only, for comparability with the Danish counterfactual of Section~\ref{sec:abm-danish}, whose simulated roll-off routinely exceeds the cap. The terminal gaps in Figure~\ref{fig:benchmark} are the cap-relative benchmark and the expectations-based complement below. One feature of the series is an artifact: four months come in at exactly zero. No result in this paper reads it as behavior. The cause is the back-out's non-negativity clip. It pins any month whose reported roll-off falls below scheduled amortization and displaces the mass into the next month --- a month-boundary allocation artifact of the reporting series (run \texttt{h1\_zero\_months\_diagnosis}). Appendix~\ref{sec:robustness-seasonalfloor} traces the mechanism and shows that repairing the affected pairs flips this paper's frozen monthly-timing rule from pass to concede. That artifact is one reason every claim made from this benchmark is a claim about levels.

\begin{figure}[H]
\centering
\includegraphics[width=0.94\textwidth]{fig9_benchmark_construction}
\caption{The empirical benchmark. Left: realized monthly SOMA agency MBS roll-off against the phased redemption cap; the shaded area is the monthly shortfall, which accrues on a net basis (the two circled opening months are net settlement inflows). Right: the same construction cumulatively. The terminal gaps are the \$764.7 billion cap-relative benchmark and, against the New York Fed's May 2022 ex-ante projection, the \$87.8 billion expectations-based complement. Monthly series from the frozen production run's metrics; the projection's monthly path is derived from the published year-end levels under the pre-committed uniform-spread rule and asserted, in the generator, to reproduce the frozen expectations-run totals.}
\label{fig:benchmark}
\end{figure}

The empirical shortfall is approximately \$764.7 billion (unrounded: \$764.748 billion; every percentage-of-benchmark figure in this paper divides by this value). It is the benchmark for simulated trapped liquidity in Sections~\ref{sec:abm} and~\ref{sec:hazard}. The SOMA face spans the Federal Reserve's full agency MBS book, Ginnie Mae collateral included, 20.4\% of face. Of that face, 66.2\% lies in the 2017--2021 vintages of the Freddie Mac loan-level sample (Section~\ref{sec:hazard}). The hazard framework therefore samples directly neither Ginnie collateral nor roughly one-third of the book by vintage. Section~\ref{sec:pathb} bounds both exclusions externally; Section~\ref{sec:limitations} states what each bound leaves open. The exclusions must be netted on the book's own joint cells, not counted separately. Out-of-window vintage is 33.7\% of face (2022 cohort 23.1\%, pre-2017 10.6\%), and the 20.4\% Ginnie share cuts across vintage. The conventional-and-in-window intersection is therefore 51.0\% on the book's own agency $\times$ vintage joint cells (bounded above by 51.8\% for cell truncation; 52.7\% under an independence approximation). The Freddie 2017--2021 conventional universe estimating Path B is about \emph{half} of SOMA book face. Two-thirds is the vintage-only marginal and overstates coverage. External-speed overlays (Section~\ref{sec:pathb}) bound Ginnie at \$20--\$47 billion by published observed speeds, and out-of-window vintages at \$11.7 billion by observed cohort speeds from the replication ingest. Each bound is signed toward overstating trapped liquidity. On the simulated side an income-scaled curtailment assumption enters, averaging approximately 0.84\% annualized and consistent with GSE daily prepayment report estimates. It is added to each simulated roll-off and netted as a common dollar item in the shared accounting layer of Section~\ref{sec:robustness-hybrid}, so the empirical CPR back-out's series retains that component. Two implementations of the back-out exist: the ABM's cohort-weighted, term-aware scheduled amortization and the hazard side's single-pool approximation. Both amortize at the pass-through WAC, though scheduled principal follows the borrowers' note rate, that WAC plus a $+0.80$-point guarantee-fee and base-servicing wedge. On the note-rate basis the hazard-side CPR level rises $+0.27$ points, a deterministic schedule difference. Every trapped-dollar figure is bit-identical: the benchmark's dollars never pass through the CPR scale. Every peak lag and $U_2$ is unchanged at reporting precision, correlations differing only at the third decimal (run \texttt{coupon\_convention\_amortization}). Table~\ref{tab:theil} is scored on the corrected basis, and Table~\ref{tab:estimators}'s note states which rows use which.

\paragraph{Accounting bases.} Every recovery percentage in this paper is quoted on one of four accounting bases, differing only in two orthogonal choices --- pure accounting layers with no modeling content --- whether the common curtailment flow is netted out, and whether the estimation sample is reweighted to the SOMA book's coupon mix. The \emph{standalone-scorer basis} makes neither adjustment; the \emph{benchmark-consistent shared basis} nets that common curtailment flow (\$69.56 billion over the window) identically from every U.S. leg, a flat 9.1 percentage points of the benchmark for every U.S. estimator; the \emph{full-book basis} instead reweights the standalone figure to the book's weighted-average coupon (WAC) without netting; and the \emph{composed basis} applies both, the same 9.1-point netting to the full-book figure. Because the netting is bit-identical across central and null legs, it cancels exactly from the lock-in marginal and every cross-estimator gap --- an invariance at every cell of the calibration box and under a $\pm 25\%$ stress on the assumed curtailment rate (derivation and stress in Section~\ref{sec:robustness-hybrid}; the four bases per hazard leg in Table~\ref{tab:bases}). The cancellation is term-by-term monthly, not merely aggregate: the netted flow is a shared monthly rate times the shared actual-holdings path, so any curtailment time profile --- calendar-seasonal, or income elasticities differing across sub-periods --- cancels identically, moving only the common wedge. The income-to-rate mapping's functional form is immaterial to the marginal for the same reason: any alternative (concave, threshold, or piecewise in income) yields another shared profile. Main-text headline figures are quoted on the shared basis throughout.

I construct a second, expectations-based benchmark because the cap-relative one carries a tension: the caps were widely expected to be non-binding for agency MBS, so a shortfall measured against them mixes an anticipated rate-inelastic component with genuine surprise. Against the Federal Reserve's own ex-ante projection, projected agency MBS runoff is \$740.6 billion and realized runoff \$652.8 billion: an expectations-based shortfall of \$87.8 billion, 11.5\% of the cap-relative \$764.7 billion. That projection is the baseline path published in the New York Fed's May 2022 Open Market Operations annual report \citep{omo2021}, starting from the February 2022 balance sheet, incorporating the FOMC's May 2022 Plans, and modeling runoff from June 2022. I spread its year-end figures uniformly within each printed year and clip them to the window under a rule fixed before the numbers were seen. The report prints year-end levels only, and no staff projection at monthly granularity is public for any vintage covering this window. Equivalently, that projection already implied 75.6--88.5\% of the realized cap-shortfall, the range spanning defensible intra-2022 allocations of its printed annual net. The pre-committed uniform-spread central construction gives \$676.9 billion, or 88.5\%. The substantive point is that a source predating every estimate in this paper anticipated the large majority of the shortfall, and it stands on its own. A sharper reading it invites does not. Under the uniform-spread construction the 88.5\% sits 0.2 points from the 88.7\% the $\beta_1 = 0$ rate-inelastic null recovers on the benchmark-consistent basis, and I presented that near-equality as convergent evidence. I withdraw the presentation: the 0.2 points require two conditions at once. They require the in-sample 4.0\% turnover floor, because at the off-window floor the paper headlines, the null recovers 85.7\% and falls 2.8 points \emph{below} the projection, not 0.2 above it. And they require the uniform-spread allocation, because the settlement-aware allocation puts the anticipated share at 75.6\%, which the null clears by 13.2 points in-sample and 10.1 points at the headline floor. A $-2.8$ to $+13.2$ point span across two disclosed choices is not a convergence, so I now report the projection range and the null's recovery separately. Both recoveries are production-floor-form reads, so the comparison is against one form's null, not a form-invariant object. Under the additive form the null recovers 35.6\% on the benchmark-consistent basis at the headline floor (44.7\% standalone; Section~\ref{sec:robustness-floor}). What survives is coarse and is what I rely on: under the production floor form both objects put the anticipated rate-inelastic component in the large majority of the shortfall. That is a fact about where the caps were set, not the size of the rate response, since the caps sat above any prepayment path this book could have delivered. Leading with it understates lock-in; the reading that speaks to mortgage markets is the marginal's share of the expectations-based shortfall, below. The central 88.5\% rides on the uniform-spread treatment of the printed 2022 annual net, which nets first-half settlement inflows against June--December runoff. That is the less faithful of the two constructions to the realized settlement record. The settlement-aware allocation assigns those inflows to the first half (realized: holdings rose \$92.3 billion over January--May 2022 in the benchmark's own SOMA series), giving roughly three-quarters (75.6\%). It leaves the large-majority reading, and every gate and threshold of the run, unchanged. Parametric within-year ramps bound the rest. Linear 2022 ramps move the share only to 88.0--89.0\%, the December-2025 window clip spans 86.1--90.6\% between all-front and linear-back extremes, and the bracket pushing all 2025 runoff to the clipped December raises it to 115.3\%. The rate-inelastic-majority threshold survives at every allocation. Every percentage-of-benchmark figure in this paper remains a cap-basis quantity. Dollar objects and estimator rankings are basis-invariant, the decomposition rescaling by the $8.7\times$ benchmark ratio on the expectations basis. The identified lock-in marginal is 49\% of the expectations-based shortfall under the central uniform-spread construction at the headline off-window floor ($+\$42.6$ billion), 80\% at the in-sample calibration point ($+\$70.3$ billion). That shortfall is itself allocation-dependent, so the ratio falls with the anticipated share (roughly 23\% and 38\% at the 75.6\% allocation) and is undefined at the extreme 115.3\% bracket, where the expectations-based shortfall turns negative. Every one of these ratios rests on one condition I cannot test: it reads the expectations-based shortfall as a residual free of lock-in, when the projection's own prepayment model may already have embedded part of the channel. That is unrecoverable from the published path of holdings. Any lock-in the projection anticipated therefore sits inside the \$740.6 billion, not the \$87.8 billion residual, and the ratios are upper bounds on lock-in's share of the genuine surprise, exact only if the projection embedded none of the channel. A second condition is likewise untestable from the published record: the projection conditions on the spring-2022 market-implied rate path, well below what was realized, so part of the residual is rate-path surprise, not lock-in. Both are properties of the projection, not of any one quotation of it, so they govern every site where these ratios appear (introduction, this section, conclusion). Each reads as an allocation-conditional upper bound on lock-in's share of a mixed lock-in-and-rate-path surprise, not a partition of it. One question of standing precedes this accounting: a shortfall against the MBS cap is not by itself evidence that a stated policy objective was missed. The Committee's operating object was the aggregate securities portfolio's decline, Treasury runoff ran against its own caps throughout, and nothing in this paper shows that reserves or the aggregate path came in off the intended course because MBS ran slow. The \$764.7 billion measures the \emph{composition} of the decline --- MBS principal that did not arrive on the book's own schedule --- and the duration consequences it carries (Section~\ref{sec:discussion-fallout}): how the decline was delivered and what that cost in portfolio structure, not a claim that the decline itself fell short.

The benchmark measures shortfall against the redemption cap, a ceiling on runoff. The cap is not a forecast. New York Fed staff projected the below-cap regime from the month the \$35 billion cap took effect \citep{nyfed2022}. The cap-relative shortfall is therefore an upper bound on the shortfall relative to any reasonable ex ante expectation. The complement above quantifies it: the expectations-based shortfall is \$87.8 billion, 11.5\% of the cap-relative figure. I retain the cap as the headline counterfactual for two reasons. First, it is the only policy-committed quantity available. The projection underlying the complement is a baseline forecast, not an FOMC-stated target. Adopting any forecast---the Federal Reserve's or a model's---would be circular with respect to models of expected prepayment, the object under study. Second, it is policy-relevant: it measures the gap between the maximum balance-sheet reduction the FOMC authorized and the reduction the mortgage structure delivered. ``Trapped liquidity'' and ``shortfall'' accordingly denote the shortfall relative to the cap throughout the paper.